\documentclass[a4paper]{article}
\usepackage{graphicx}
\usepackage{float}
\usepackage{wrapfig}
\usepackage[a4paper, margin=2.5cm]{geometry}

% --- Encoding & Language ---
\usepackage[utf8]{inputenc}     % pdfLaTeX
\usepackage[T1]{fontenc}
\usepackage[magyar]{babel}

% --- Math & Chemistry ---
\usepackage{amsmath, amsthm, amssymb}
\usepackage[version=4]{mhchem} % chemical formulas
\usepackage{siunitx}            % units, temperatures
\usepackage{tikz}
\usetikzlibrary{tikzmark, calc}
\usepackage{mathtools}

% --- Section formatting ---
\usepackage{titlesec}

% Allow deep section numbering and TOC
\setcounter{secnumdepth}{5}
\setcounter{tocdepth}{5}

\makeatletter

% ------------------------
% LEVEL 4: subsubsubsection
% ------------------------
\titleclass{\subsubsubsection}{straight}[\subsubsection]
\newcounter{subsubsubsection}[subsubsection]
\renewcommand\thesubsubsubsection{\thesubsubsection.\arabic{subsubsubsection}}

\renewcommand\subsubsubsection{%
  \@startsection{subsubsubsection}{4}{\z@}%
  {3.25ex \@plus1ex \@minus.2ex}%
  {1.5ex \@plus .2ex}%
  {\normalfont\normalsize\bfseries}%
}

\titleformat{\subsubsubsection}
  {\normalfont\normalsize\bfseries}
  {\thesubsubsubsection}{1em}{}

\titlespacing*{\subsubsubsection}{0pt}{3ex}{1ex}

% TOC entry for level 4
\newcommand*\l@subsubsubsection{\@dottedtocline{4}{7em}{4em}}

% --------------------------
% LEVEL 5: subsubsubsubsection
% --------------------------
\titleclass{\subsubsubsubsection}{straight}[\subsubsubsection]
\newcounter{subsubsubsubsection}[subsubsubsection]
\renewcommand\thesubsubsubsubsection{\thesubsubsubsection.\arabic{subsubsubsubsection}}

\renewcommand\subsubsubsubsection{%
  \@startsection{subsubsubsubsection}{5}{\z@}%
  {3.25ex \@plus1ex \@minus.2ex}%
  {1.5ex \@plus .2ex}%
  {\normalfont\normalsize\bfseries}%
}

\titleformat{\subsubsubsubsection}
  {\normalfont\normalsize\bfseries}
  {\thesubsubsubsubsection}{1em}{}

\titlespacing*{\subsubsubsubsection}{0pt}{3ex}{1ex}

% TOC entry for level 5
% Use level 4 internally for safety (LaTeX only knows TOC levels 1-4)
\newcommand*\l@subsubsubsubsection{\@dottedtocline{4}{10em}{5em}}

\makeatother

\title{IV. Főcsoport: A széncsoport}
\author{Rovenszky Robin}
\date{2026. Május 19.}

\newtheorem{definition}{Definíció}

\begin{document}

\maketitle
\tableofcontents

\vspace{14pt}
\hrule

\section{A széncsoport}

4 db vegyérték elektron\\
$\rightarrow$ 4 db kovalens kötés (C, Si)

\noindent
\begin{minipage}{0.78\textwidth}
\begin{enumerate}
    \item Szén (C)
    \item Szilícium (Si)
    \item Germánium (Ge)
    \item Ón (Sn)
    \item Ólom (Pb)
\end{enumerate}

5 db vegyérték elektron $\implies$ a nemesgáz elektronszerkezetének eléréséhez:
    \begin{itemize}
        \item 3 db kovalens kötés\\
        Pl.: Ammónia 
        \vspace{-10pt}
        \begin{figure}[H]
            \centering
            \includegraphics[width=0.3\textwidth]{Assets/Chemistry/Main Group V/NH_3.png}
        \end{figure}
        \vspace{-30pt}
    \end{itemize}
    
    \end{minipage}
    \hfill
    \begin{minipage}{0.1\textwidth}
        \centering
        
\includegraphics[width=\linewidth]{Assets/Chemistry/Main Group IV/Periodic_Table.png}
    \end{minipage}
        
    \begin{itemize}
    
        \item 3-szoros negatív töltésű anion (pl. $N^{3-}$ - nitrid-ion). \\
        Egyedi, mert lényegében csak a nitrogén tud ennyi iont felvenni, mivel nagyon nagy az elektronegativitása.
        \item Az is lehetséges, hogy a nitrogén 5 db kovalens kötést létesítsen. Pl.: $HNO_3$ - salétromsav

        \begin{figure}[H]
            \centering
            \includegraphics[width=0.3\textwidth]{Assets/Chemistry/Main Group V/HNO_3.png}
        \end{figure}
        
        A nitrogén túllépett a nemesgáz elektroszerkezetén, de ez így is energetikailag kedvező.
    \end{itemize}

\begin{table}[h!]
    \centering
    \begin{tabular}{p{2.5cm}|p{4cm}|p{4cm}}
    
    Szempont & $CO_2$ & $CO$ \\ 
    \hline
    
    Keletkezés &
    Szén égésekor:
    $C + O_2 \rightarrow CO_2$

    C-tartalmú anyag égésekor.
    Pl:
    $CH_4 + 2O_2 \rightarrow CO_2 + 2H_2O$
    & 
    a\\
    
    \hline
    
    Élettani hatás & 
    a & 
    a\\
    
    \hline

    Előfordulás & 
    a & 
    a \\
    
    \hline
    
    Környezeti hatás & 
    a &
    a \\
    
    \hline
    
    Felhasználás & 
    a & 
    a
    
    \end{tabular}
    
    \caption{Szén vegyületek}
    \label{tab:example}
\end{table}

\end{document}